\chapter{event log：如何验证订单因果链}

\section{本章要解决的困惑}

fast 回测的输出是表格。strict replay 的核心输出是 compact event log：

\begin{lstlisting}
runs/<run_id>/events.ndjson
\end{lstlisting}

它不是为了好看，而是为了回答：

\begin{quote}
这笔订单为什么出现？什么时候到达？按什么价格成交？成交后账户如何变化？
\end{quote}

\section{NDJSON 为什么适合 event log}

NDJSON 是一行一个 JSON 事件。优点是：

\begin{itemize}[leftmargin=2em]
  \item 可以流式写入，不需要把整天事件放内存；
  \item 崩溃时前面的事件仍然可读；
  \item Monitor 和 audit 工具可以逐行解析；
  \item compact log 可以只记录 causality，不记录每帧 hold。
\end{itemize}

当前 compact log 不应该记录每个市场事件上的 hold，因为那会把日志膨胀成低价值流水账。

\section{事件 envelope}

\code{event_contracts.md} 中的 run event log envelope：

\begin{lstlisting}
event_id
run_id
event_type
replay_seq
replay_ts
wall_ts_ms
source
payload
\end{lstlisting}

读事件时，先看 \code{event_type}，再看 \code{payload}。

\section{订单因果链}

一笔严格 replay 的链条大致是：

\begin{lstlisting}
market event
-> strategy_feature_checkpoint
-> order_intent
-> order_scheduled
-> order_arrived
-> order_ack or order_reject
-> fill_created
-> portfolio_state_on_change
\end{lstlisting}

Monitor 的订单链表就是把这些事件按 \code{intent_id} 或相关 id 聚合起来。它不创造事实，只把事实排成可读链。

\section{一个因果链要检查什么}

\begin{longtable}{p{0.26\textwidth}p{0.64\textwidth}}
\toprule
检查项 & 问题 \\
\midrule
\code{observed_seq} & Bot 是看到哪一个市场事件后下单的？ \\
\code{arrival_seq} & 订单到达时 Runner 选的是哪个 quote/book？ \\
\code{side} & buy 还是 sell，是否和 TFI 方向一致？ \\
\code{fill_price} & buy 是否打 arrival ask，sell 是否打 arrival bid？ \\
\code{arrival_quote_lag} & 到达 quote 是否滞后？ \\
\code{latency_slippage_bps} & observed 到 arrival 之间损失多少？ \\
\code{portfolio_state} & 成交后 position/equity 是否更新？ \\
\bottomrule
\end{longtable}

\section{event log 和 entries.parquet 的关系}

\code{entries.parquet} 是 fast/research 表格输出。event log 是 strict replay 的因果审计输出。

它们回答的问题不同：

\begin{itemize}[leftmargin=2em]
  \item entries/exits：这个策略在 fast 近似下表现如何？
  \item events.ndjson：这个订单在本地交易所里到底发生了什么？
\end{itemize}

不要用 event log 重新发明策略，也不要让 Monitor 从日志里创造没有发生过的订单。

\section{手动检查点}

看到一笔 \code{fill_created}，往前找：

\begin{lstlisting}
same intent/order id
-> order_ack
-> order_arrived
-> order_scheduled
-> order_intent
-> strategy_feature_checkpoint
\end{lstlisting}

如果往前找不到完整链，先做 audit，不要先讨论策略收益。

